\section{Conclusion}
\label{sec:conclusion}

This paper has taken the view that the right way to evaluate mobile GUI agents is at the level of cross-app and cross-category generalization, and that the right mechanism for the deployment-time budget is to co-adapt a structured in-context knowledge representation alongside the policy's weights. The two-layer FSM gives ``transferable'' a structural meaning enforced by a linter rather than left to stylistic discipline. The EvoFSM-RL loop ties the upper layer's online evolution to a group-relative weight update, runs identically at source-pool pretraining and at target-app deployment, and produces measurable lifts on near-transfer apps---$+9.3$~pp from static category knowledge, $+3.7$~pp from online evolution---with a clean null control on far-transfer that confirms the gains are mechanism-bound and not artifact-bound.

The framework is not a silver bullet. Joint test-time adaptation lifts performance where the base policy has nonzero success on the adaptation set, and degrades gracefully to no update when the base policy fails uniformly. We see this both as a real limitation and as a precise characterization of where the next progress has to come from: reward shaping that creates signal on apps the base policy cannot already complete in part, or stronger base policies on which the same joint loop has more to work with. Either direction inherits the structural pieces this paper contributes---the two-layer separation, the category-level abstract library, and the joint pretraining-then-deployment loop---and can be evaluated under the same strict pool, category, and template disjoint protocol.

More broadly, the design space between hand-written agent prompts and full per-app fine-tuning is wider than it has been treated. Structuring an agent's prompt so that one layer of it can be safely mutated by an automated, RL-coupled process is a small commitment with a clear payoff. The two-layer FSM is one concrete instantiation, and we expect similar recipes to admit other structured forms---state-and-skill libraries, retrieved code snippets, executable verifiers---in agent settings where the underlying task vocabulary is rich enough to support such structure.
